% =========================================================
\chapter{Les essais mécaniques}
% =========================================================

\begin{rembox}[title=En bref]
\textbf{De quoi parle ce chapitre ?} Avant de parler des matériaux, il faut savoir
\textbf{comment on mesure} leur comportement mécanique. Le chapitre passe en revue les essais
classiques : traction (le roi), compression, dureté, flexion, essais où le temps compte
(vitesse de déformation, relaxation, fluage), puis les essais de rupture (ténacité, Charpy,
formabilité, fatigue). \
\textbf{Les trois idées à retenir} : la différence entre grandeurs \emph{nominales} et
\emph{vraies} ; la \textbf{striction} comme instabilité (critère de Considère) ; la
\textbf{mécanique de la rupture} ($K = Y\sigma\sqrt{\pi a}$), qui reviendra dans tous les
chapitres suivants.
\end{rembox}

\section{Pourquoi l'essai de traction est le roi des essais}

\slide[0.75]{c01-06}{Le barreau de traction : une section réduite (\emph{gauge section}) entre deux zones d'amarrage.}

On tire sur un barreau dont le milieu est aminci. Tout l'intérêt est là : dans la section
réduite, l'état de contrainte est le plus simple possible --- \textbf{une seule composante,
$\sigma_{xx}$, la même partout}. Toutes les autres composantes du tenseur des contraintes sont
nulles. On sait donc exactement ce qu'on impose au matériau, ce qui n'est vrai pour presque
aucun autre essai. Les zones d'amarrage sont plus larges pour que ce soit la section réduite,
et elle seule, qui se déforme et casse.

La déformation, elle, a \textbf{trois} composantes : le barreau s'allonge selon $x$ mais
rétrécit selon $y$ et $z$.

\begin{thmbox}[title=Élasticité linéaire isotrope en traction]
\[ \varepsilon_{xx} = \frac{\sigma_{xx}}{E}, \qquad
   \varepsilon_{yy} = \varepsilon_{zz} = -\nu\,\frac{\sigma_{xx}}{E} \]
$E$ = module d'Young (rigidité), $\nu$ = coefficient de Poisson (combien le matériau se
rétrécit latéralement quand on l'allonge).
\end{thmbox}

\begin{propbox}[title={À volume constant, $\nu = 1/2$}]
Le changement relatif de volume vaut, pour de petites déformations,
$\dfrac{\Delta V}{V} \approx \varepsilon_{xx} + \varepsilon_{yy} + \varepsilon_{zz}$.
Si le volume ne change pas (c'est \textbf{grosso modo le cas en déformation plastique}),
cette somme est nulle, donc $\varepsilon_{yy} = \varepsilon_{zz} = -\varepsilon_{xx}/2$ :
tout se passe comme si $\nu = 0{,}5$. En élasticité, $\nu$ est plus petit (souvent
$\approx 0{,}3$ pour les métaux, $\approx 0{,}25$ pour le verre) : le volume augmente un peu
quand on tire.
\end{propbox}

\begin{rembox}[title=En pratique]
La force $P$ est mesurée par une \textbf{cellule de charge}. L'allongement $\Delta L$ se lit
sur le déplacement de la traverse (peu précis, car toute la machine se déforme aussi), ou mieux
avec un \textbf{extensomètre} pincé sur la section réduite, une \textbf{jauge résistive} collée
dessus, ou par suivi optique de points à la surface.
\end{rembox}

\begin{goatbox}
Tirer sur un barreau, c'est la seule façon simple d'être sûr de ce que l'on fait subir au
matériau : partout au milieu du barreau, il est tiré de la même façon, dans une seule
direction. Quand il s'allonge, il maigrit sur les côtés, comme un chewing-gum qu'on étire.
\end{goatbox}
